% ─────────────────────────────────────────────────────────────────────────
%  vimstyle.tex — extra preamble for the Vim/Neovim note
%  (keycaps, key tables, source-code listings)
% ─────────────────────────────────────────────────────────────────────────
\usepackage{tabularx}
\usepackage{longtable}
\usepackage{listings}
\usepackage{upquote}

% ── Inline notation ──────────────────────────────────────────────────────
% \key{dw}   a keystroke, drawn as a small keycap
% \cmd{:w}   an Ex command / option name
% \file{~/.config/nvim/init.lua}
\newcommand{\key}[1]{%
	{\setlength{\fboxsep}{2pt}%
		\fcolorbox{Gray!45}{Gray!8}{\small\ttfamily #1}}}
\newcommand{\cmd}[1]{\texttt{\color{MidnightBlue!85!black}#1}}
\newcommand{\file}[1]{\texttt{\color{RawSienna!85!black}#1}}

% ── Key tables ───────────────────────────────────────────────────────────
% \begin{keys} \key{dw} & delete a word \\ ... \end{keys}
\newenvironment{keys}[1][0.26]{%
	\par\vspace{0.35em}\noindent
	\renewcommand{\arraystretch}{1.25}%
	\tabularx{\linewidth}{@{}>{\raggedright\arraybackslash}p{#1\linewidth}>{\raggedright\arraybackslash}X@{}}%
}{%
	\endtabularx\par\vspace{0.35em}%
}

% Same, but for the long reference tables in the appendix (page breakable)
\newenvironment{keyslong}[1][0.26]{%
	\renewcommand{\arraystretch}{1.25}%
	\begin{longtable}{@{}>{\raggedright\arraybackslash}p{#1\linewidth}>{\raggedright\arraybackslash}p{\dimexpr0.96\linewidth-#1\linewidth\relax}@{}}%
		}{%
	\end{longtable}%
}

% ── Listings ─────────────────────────────────────────────────────────────
\lstdefinestyle{nvimlua}{
	language=[5.2]Lua,
	basicstyle=\footnotesize\ttfamily,
	keywordstyle=\color{MidnightBlue}\bfseries,
	commentstyle=\color{Gray!75!black}\itshape,
	stringstyle=\color{RawSienna},
	identifierstyle=\color{black},
	showstringspaces=false,
	columns=fullflexible,
	keepspaces=true,
	breaklines=true,
	breakatwhitespace=true,
	tabsize=2,
	morekeywords={vim,require,pairs,ipairs,setmetatable,opts,true,false,nil},
	frame=none,
	xleftmargin=0pt,
	xrightmargin=0pt,
	aboveskip=0pt,
	belowskip=0pt,
}

\lstdefinestyle{vimsession}{
	basicstyle=\footnotesize\ttfamily,
	columns=fullflexible,
	keepspaces=true,
	breaklines=true,
	tabsize=4,
	commentstyle=\color{Gray!75!black}\itshape,
	morecomment=[l]{\#},
	frame=none,
	xleftmargin=0pt,
	xrightmargin=0pt,
	aboveskip=0pt,
	belowskip=0pt,
}

% Code blocks are drawn by tcolorbox rather than by listings' own `frame'.
% listings computes its frame from \linewidth, which is wrong inside the
% theorem boxes of header.tex: the top rule comes out shifted and overhanging.
% tcolorbox gets the geometry right in every context, so the styles above set
% frame=none and the box below supplies the rule and the padding.
%
% These take no optional argument on purpose: with one, LaTeX looks ahead past
% the newline for a `[', and a listing whose first line is a bare `{' (a Lua
% table, say) is then swallowed as a group instead of being typeset.
\tcbuselibrary{listings}
\tcbset{
	code box/.style={
		enhanced, sharp corners,
		colframe=Gray!45, colback=Gray!5, boxrule=0.5pt,
		left=7pt, right=7pt, top=5pt, bottom=5pt,
		before skip=1em, after skip=1em,
		splitsafe,                % see header.tex: breakable + colback is unsafe
	                              % under XeTeX, so these stay whole instead
	}
}

\newtcblisting{luacode}{code box,
	listing only, listing options={style=nvimlua}}
\newtcblisting{vimcode}{code box,
	listing only, listing options={style=vimsession}}
\newtcblisting{shell}{code box,
	listing only, listing options={style=shellsession}}
\newtcblisting{bashcode}{code box,
	listing only, listing options={style=bashscript}}

% ── Makefile listings ────────────────────────────────────────────────────
% `makecode' — a Makefile excerpt.  Recipe indentation is written as spaces
% here, because tcolorbox's verbatim reader does not pass tabs through to
% listings; sec:mk-notation says so in the text, and sec:mk-tab explains why a
% real Makefile needs the tab.
\lstdefinestyle{makefile}{
	language=[gnu]make,
	basicstyle=\footnotesize\ttfamily,
	keywordstyle=\color{MidnightBlue}\bfseries,
	commentstyle=\color{Gray!75!black}\itshape,
	stringstyle=\color{RawSienna},
	showstringspaces=false,
	columns=fullflexible,
	keepspaces=true,
	breaklines=true,
	tabsize=4,
	frame=none,
	xleftmargin=0pt,
	xrightmargin=0pt,
	aboveskip=0pt,
	belowskip=0pt,
}


\newtcblisting{makecode}{code box,
	listing only, listing options={style=makefile}}

% ── C and GDB listings ───────────────────────────────────────────────────
% `ccode'    — a C source excerpt.
% `gdbcode'  — a GDB session: "(gdb)" prompt lines are input, the rest output.
\lstdefinestyle{csource}{
	language=C,
	basicstyle=\footnotesize\ttfamily,
	keywordstyle=\color{MidnightBlue}\bfseries,
	commentstyle=\color{Gray!75!black}\itshape,
	stringstyle=\color{RawSienna},
	showstringspaces=false,
	columns=fullflexible,
	keepspaces=true,
	breaklines=true,
	tabsize=4,
	numbers=left,
	numberstyle=\scriptsize\color{Gray!90!black},
	numbersep=9pt,
	frame=none,
	xleftmargin=16pt,   % room for the line numbers inside the box
	xrightmargin=0pt,
	aboveskip=0pt,
	belowskip=0pt,
}

\lstdefinestyle{gdbsession}{
	basicstyle=\footnotesize\ttfamily,
	columns=fullflexible,
	keepspaces=true,
	breaklines=true,
	tabsize=4,
	commentstyle=\color{Gray!75!black}\itshape,
	morecomment=[l]{\#},
	frame=none,
	xleftmargin=0pt,
	xrightmargin=0pt,
	aboveskip=0pt,
	belowskip=0pt,
}

\newtcblisting{ccode}{code box,
	listing only, listing options={style=csource}}
\newtcblisting{gdbcode}{code box,
	listing only, listing options={style=gdbsession}}

% `conf' — a configuration file: .gitconfig, .gitignore, YAML, INI.
\lstdefinestyle{conffile}{
	basicstyle=\footnotesize\ttfamily,
	% Latin Modern Mono ligates << and >> into guillemets, which mangles Git
	% conflict markers.  Emitting each character separately suppresses it.
	literate={<}{{\char60}}1 {>}{{\char62}}1,
	columns=fullflexible,
	keepspaces=true,
	breaklines=true,
	tabsize=4,
	commentstyle=\color{Gray!75!black}\itshape,
	morecomment=[l]{\#},
	stringstyle=\color{RawSienna},
	frame=none,
	xleftmargin=0pt,
	xrightmargin=0pt,
	aboveskip=0pt,
	belowskip=0pt,
}

\newtcblisting{conf}{code box,
	listing only, listing options={style=conffile}}


% ── A "try it" box, for hands-on drills ──────────────────────────────────
\newenvironment{tryit}{%
	\begin{tcolorbox}[
			enhanced, sharp corners, arc=4pt, outer arc=4pt,
			colframe=SeaGreen!70!black, colback=SeaGreen!3, boxrule=0.6pt,
			before skip=10pt, after skip=14pt, splitsafe,
			left=10pt, right=10pt, top=5pt, bottom=6pt,
			title={\sffamily\bfseries\small Try it},
			coltitle=white, colbacktitle=SeaGreen!65!black,
			fonttitle=\sffamily\bfseries\small]%
		}{%
	\end{tcolorbox}%
}

% ── Cross-reference names ────────────────────────────────────────────────
% The template maps \chapterautorefname to "Section", which is fine for a
% single-subject note but ambiguous here: this book has both Chapter 3 and
% Section 3.1.  Name them separately.
\def\chapterautorefname{Chapter}
\def\sectionautorefname{Section}
\def\subsectionautorefname{Section}
\def\partautorefname{Part}

% ── Shell listings ───────────────────────────────────────────────────────
% `shell'     — a terminal transcript: prompt lines plus program output.
% `bashcode'  — a script file, with shell keywords highlighted.
\lstdefinestyle{shellsession}{
	basicstyle=\footnotesize\ttfamily,
	columns=fullflexible,
	keepspaces=true,
	breaklines=true,
	breakatwhitespace=false,
	tabsize=4,
	commentstyle=\color{Gray!75!black}\itshape,
	morecomment=[l]{\#},
	frame=none,
	xleftmargin=0pt,
	xrightmargin=0pt,
	aboveskip=0pt,
	belowskip=0pt,
}

\lstdefinestyle{bashscript}{
	style=shellsession,
	language=sh,
	keywordstyle=\color{MidnightBlue}\bfseries,
	morekeywords={local,return,function,source,export,declare,readonly,shift,
		trap,set,unset,printf,echo,read,exit,eval,shopt,select},
	stringstyle=\color{RawSienna},
	showstringspaces=false,
	alsoletter={-},
}



% Three-digit page numbers need a wider column in the table of contents.
\DeclareTOCStyleEntries[pagenumberwidth=2.2em]{tocline}{part,chapter,section}
